\chapter{مقدمه‌ای بر داده‌های پیچیده}
\label{ch:complex-data}

\section*{اهداف این فصل}
\begin{itemize}
    \item آشنایی مقدماتی با داده‌کاوی گراف
    \item آشنایی مقدماتی با تحلیل سری‌های زمانی
    \item آشنایی با مسیرهای پیشرفته‌تر برای ادامهٔ یادگیری
\end{itemize}

\section{فراتر از داده‌های جدولی}

در تمام فصل‌های پیشین، فرض بر این بود که داده به‌صورت جدولی است: هر
سطر یک نمونهٔ مستقل و هر ستون یک ویژگی. بسیاری از داده‌های واقعی، اما،
ساختاری فراتر از این دارند — روابط میان نمونه‌ها (مانند شبکه‌های
اجتماعی) یا ترتیب زمانی میان مشاهدات (مانند قیمت سهام). این فصل پایانی،
به‌صورت مقدماتی و صرفاً برای آشنایی، دو نوع از این داده‌های پیچیده‌تر
را معرفی می‌کند تا خواننده افق‌های بعدی مطالعه را بشناسد.

\section{داده‌کاوی گراف}

در بسیاری از مسائل، داده به‌طور طبیعی به‌صورت
\dmterm{\textbf{گراف}}{Graph} — مجموعه‌ای از
\dmterm{گره‌ها}{Nodes / Vertices} و
\dmterm{یال‌ها}{Edges} میان آن‌ها — بازنمایی می‌شود: کاربران یک شبکهٔ
اجتماعی و روابط دوستی میانشان، صفحات وب و پیوندهای میانشان، یا مولکول‌ها
و پیوندهای شیمیایی. \dmterm{\textbf{داده‌کاوی گراف}}{Graph Mining} به
کشف الگو در این نوع داده می‌پردازد.

برخی از سؤالات رایج در این حوزه عبارت‌اند از:

\begin{itemize}
    \item \dmterm{\textbf{تشخیص اجتماع}}{Community Detection}: یافتن گروه‌هایی از گره‌ها
          که درون خود پیوندهای متراکم‌تری نسبت به بیرون دارند — مفهومی
          که شباهت زیادی به خوشه‌بندی (فصل \ref{ch:clustering}) دارد،
          با این تفاوت که ورودی آن روابط میان نمونه‌هاست، نه ویژگی‌های
          عددی هر نمونه.
    \item \dmterm{\textbf{مرکزیت گره}}{Centrality}: سنجش میزان اهمیت یا نفوذ هر گره در
          گراف (مثلاً شناسایی کاربران تأثیرگذار در یک شبکهٔ اجتماعی).
    \item \dmterm{\textbf{پیش‌بینی پیوند}}{Link Prediction}: پیش‌بینی این‌که کدام یال‌های
          جدید در آینده به گراف اضافه خواهند شد (مثلاً پیشنهاد دوستی در
          شبکه‌های اجتماعی).
\end{itemize}

\section{تحلیل سری‌های زمانی}

\dmterm{\textbf{سری زمانی}}{Time Series} دنباله‌ای از مشاهدات است که
به ترتیب زمانی ثبت شده‌اند (مثلاً دمای روزانه، قیمت سهام، یا میزان
فروش ماهانه). تفاوت اساسی این نوع داده با داده‌های جدولی معمولی، وجود
\textbf{وابستگی زمانی} میان مشاهدات متوالی است؛ برخلاف فرض رایج در
بسیاری از الگوریتم‌های پیشین این کتاب که نمونه‌ها را مستقل از هم در
نظر می‌گرفتند.

\dmterm{\textbf{تجزیهٔ سری زمانی}}{Time Series Decomposition} یک سری
زمانی را معمولاً به سه مؤلفهٔ اصلی تجزیه می‌کند:

\begin{itemize}
    \item \dmterm{\textbf{روند}}{Trend}: تغییر کلی و بلندمدت (صعودی یا
          نزولی) در داده.
    \item \dmterm{\textbf{فصلی بودن}}{Seasonality}: الگوهای تکرارشونده
          در بازه‌های زمانی ثابت (مثلاً افزایش فروش هر سال در نزدیکی
          یک مناسبت خاص).
    \item \dmterm{\textbf{باقیمانده}}{Residual}: نوسانات تصادفی باقی‌مانده
          پس از کنار گذاشتن روند و فصلی بودن.
\end{itemize}

برای پیش‌بینی مقادیر آیندهٔ یک سری زمانی، از مدل‌های کلاسیکی مانند
\dmterm{\LR{ARIMA}}{AutoRegressive Integrated Moving Average} تا
روش‌های مدرن مبتنی بر یادگیری عمیق (مانند شبکه‌های بازگشتی\footnote{\LR{Recurrent Neural Networks (RNN)}})
استفاده می‌شود که بررسی
کامل آن‌ها خارج از حوصلهٔ این کتاب مقدماتی است.

\section{جمع‌بندی و مسیرهای پیشرفته‌تر}

در طول این کتاب، مسیری از مفاهیم پایه (فصل ۱ و ۲)، از میان یادگیری با
نظارت (طبقه‌بندی و رگرسیون در فصل‌های ۳ تا ۵)، تا یادگیری بدون نظارت
(قوانین انجمنی، خوشه‌بندی و تشخیص داده پرت در فصل‌های ۶ تا ۸)، و در
پایان مقدمه‌ای بر داده‌های متنی و پیچیده (فصل‌های ۹ و ۱۰) طی کردیم. این
کتاب صرفاً نقطهٔ شروع است؛ برای ادامهٔ مسیر، خواننده می‌تواند به سراغ
موضوعاتی مانند یادگیری عمیق، یادگیری تقویتی، پردازش زبان طبیعی مدرن (بر
پایهٔ ترنسفورمرها)، و مقیاس‌پذیری داده‌کاوی روی داده‌های عظیم (با
ابزارهایی مانند Spark) برود. مراجع اصلی این کتاب — Han/Kamber/Pei و
Tan/Steinbach/Kumar — نیز برای مطالعهٔ عمیق‌تر هر یک از فصل‌های این
کتاب به‌شدت توصیه می‌شوند.

\section*{تمرین‌ها}

\begin{enumerate}
    \item شباهت و تفاوت «تشخیص اجتماع» در گراف با «خوشه‌بندی» در
          فصل \ref{ch:clustering} را توضیح دهید.
    \item یک مثال واقعی از پیش‌بینی پیوند در یک شبکهٔ اجتماعی که با آن
          آشنایی دارید، بیان کنید.
    \item برای فروش ماهانهٔ یک فروشگاه لباس، مؤلفه‌های روند، فصلی بودن
          و باقیمانده را با یک مثال فرضی توضیح دهید.
    \item چرا فرض استقلال نمونه‌ها که در بسیاری از الگوریتم‌های فصل‌های
          قبل وجود داشت، برای داده‌های سری زمانی مناسب نیست؟
    \item با توجه به علایق خودتان، یکی از مسیرهای پیشرفته‌تر معرفی‌شده
          در این فصل (یادگیری عمیق، یادگیری تقویتی، پردازش زبان طبیعی
          مدرن، یا داده‌کاوی مقیاس‌بزرگ) را انتخاب کنید و توضیح دهید
          چرا برای ادامهٔ مطالعه به آن علاقه‌مندید.
\end{enumerate}
